\chapter*{译者说明}
\addcontentsline{toc}{chapter}{译者说明}
\markboth{译者说明}{译者说明}% 否则续页页眉还留着上一章（目录）的记号

\section*{底本与范围}
本译本以 Laurent Fargues 与 Peter Scholze 的《Geometrization of the local Langlands
correspondence》（arXiv:2102.13459v4，2024 年 11 月 27 日版，全书 356 页）为底本。
底本的印刷页码与 PDF 页码一致（偏移为 0），因此本书所有溯源标记都直接写原文页码。

\textbf{当前进度}：第一章（导论，原文 p.7–44）已完整译出并校订；其余各章按章推进，
每章完成即更新本译本。未完成的章在目录中暂不出现，不以空章占位。

\section*{逐页溯源}
正文左侧空白处的〔p.$N$〕是\textbf{原文页码锚}，表示从此处起对应原文第 $N$ 页。
遇到页界正好落在句子中间、边注排不下的地方，改用行内小上标〔p.$N$〕。
两类标记都会同时写入编译台账，可用 \texttt{工程文件/tools/check\_pagemarks.py}
核查「有没有真的印出来」，用 \texttt{page\_map.py} 生成原文页→译本页对照表。

\section*{译注的三类来源}
本书的解释性注释一律用楷体小字排，并强制标注来源，分三级：

\begin{itemize}
  \item \textbf{A}：原文自身明写（含前后文、原文脚注、原文表格），只做定位与串联；
  \item \textbf{B}：本库既有译本与资料卡（如 \texttt{datasets/} 下的既有工程）中的对应表述；
  \item \textbf{I}：我们的推断或补足。若连一条外部依据都找不到，会明确写成
        「译者按（无外部依据）」，不冒充有据。
\end{itemize}

\textbf{不作}的事：不替作者补写证明，不把「作者说可能是」译成「就是」，不隐去原文的
猜想语气；原文的编号、公式、记号一律照录，只在确证笔误时改并在勘误表登记。

\section*{关于脚注编号}
本书的\textbf{译注}与作者的\textbf{原文脚注}共用同一条脚注编号序列（都排在页脚），
所以原文脚注的编号在本译本里不一定连续——例如第一章原文注 1--4 在本译本中显示为
4、6、10、17 号。引述原文脚注时请以它所在的〔p.$N$〕定位，不要按编号去找。
凡译注都在句末标出「来源：A/B/I」，没有这行标注的脚注即作者原注。

\section*{术语与勘误}
全书术语在开工前冻结（\texttt{工程文件/work/glossary.tsv} 与
\texttt{术语裁定.md}），跨章统一：perfectoid 作「完满」、diamond 作「菱形」、
condensed 作「凝聚态」、solid 作「固体」、excursion operator 作「巡游算子」、
Langlands dual group 作「朗兰兹对偶群」、Hecke eigensheaf 作「Hecke 特征层」；
shtuka、$\Bun_G$、anima、ULA 等不译。

译校过程中发现的原文笔误登记在 \texttt{工程文件/work/errata.tsv}（含位置、原写法、
订正、两处独立依据、置信度、影响卡片、状态）。置信度不足两处独立依据的只记为
「存疑」，正文保留原文写法，不擅自改。

\section*{权利与用途}
译文与语料是对第三方底本的演绎作品，本工程不代原作者授权、也不另设限制；
权利人通知即撤。自产的术语表、工程文档、工作流与资料卡按 CC BY 4.0 发布。
